\documentclass[UTF8,a4paper,11pt]{ctexart}
\usepackage{amsmath,amssymb}
\usepackage[margin=2.5cm]{geometry}
\usepackage[hidelinks]{hyperref}
\title{2014 年安徽高考文理共同题\\三次函数与边界参数质检}
\date{2026 年 10 月 24 日}
\begin{document}
\maketitle
\noindent\textbf{出处：}文科第 20 题 \url{https://github.com/DxAThing/Gaokao-Math-Problems-Compilation/blob/f1ae169b58a75b3db6e35c670b3f0d5c78718188/content/2014/anhui_liberal.tex#L628}

理科第 18 题 \url{https://github.com/DxAThing/Gaokao-Math-Problems-Compilation/blob/f1ae169b58a75b3db6e35c670b3f0d5c78718188/content/2014/anhui_science.tex#L279}

设 $f_a(x)=1+(1+a)x-x^2-x^3$，其中 $a>0$。
\paragraph{全实数的单调性}
\[
 f_a'(x)=1+a-2x-3x^2=-3(x-\alpha)(x-\beta)
\]
其中
\[
 \alpha=\frac{-1-\sqrt{4+3a}}3<0,\quad
 \beta=\frac{-1+\sqrt{4+3a}}3>0.
\]
导数符号由左至右为负、正、负。
所以函数在 $(-\infty,\alpha)$ 和
$(\beta,+\infty)$ 严格递减，在 $(\alpha,\beta)$ 严格递增。

\paragraph{$[0,1]$ 上最大值的取点}
由 $\beta>0$ 及
$\beta\le1\Longleftrightarrow a\le4$，可得
\[
 x_{\max}=
 \begin{cases}
  \dfrac{-1+\sqrt{4+3a}}3,&0<a<4,\\[2mm]
  1,&a\ge4.
 \end{cases}
\]

\paragraph{$[0,1]$ 上最小值的取点}
函数在该区间内先增后减或一直增，
内点不可能取得小于两端点的函数值。
比较 $f_a(0)=1$、$f_a(1)=a$，得到
\[
 x_{\min}=
 \begin{cases}
  1,&0<a<1,\\
  0\ \text{或}\ 1,&a=1,\\
  0,&a>1.
 \end{cases}
\]
特别地，$a=1$ 两个端点均取得最小值。

\paragraph{上游质检发现}
文科 LaTeX 的解析结尾缺少上式的完整分类，
理科版本解析完整。本稿保留两份原始卷别出处并独立补正。

\noindent\textbf{核验边界：}数学证明已独立复核；
新 TeX 尚未实际编译。
\end{document}
